%------------------------------------------------------------------------------%

\usepackage{calculo_varias_variaveis-1}

\title {Diferenciabilidade}

%------------------------------------------------------------------------------%
\begin{document}

\addtitle

%------------------------------------------------------------------------------%
\section{Diferenciabilidade}

%------------------------------------------------------------------------------%
\begin{frame}{Derivada Ordinária -- Uma variável}

  Variação da função com relação a variável

  \pause
  \vspace{\baselineskip}
  Derivada da função \(f\colon R\subset\R\to\R\) no ponto $x_0$

  \[
    \frac{df}{dx}\evalat{x=x_0}{}
    =\; \lim_{h\to 0} \frac{f(x_0+h)-f(x_0)}{h}
  \]
\end{frame}

%------------------------------------------------------------------------------%
\begin{frame}{Diferenciabilidade em 1D}
	Se a \emph{derivada de $f$ existe} em $x_0$ então $f$ é \emph{diferenciável} em $x_0$

	\pause
	\vspace{\baselineskip}
	Existe uma \emph{reta tangênte} a $f$ no ponto \(\left(x_0, f(x_0)\right)\)

	\pause
	\vspace{\baselineskip}
	Podemos \emph{aproximar} a função pela reta em uma vizinhança de $x_0$
	\[
        \pause
        \Delta y = f(x) - f(x_0)
    \]
    \[
	 	\pause
	 	\Delta y = f'(x_0)\Delta x + \varepsilon \Delta x
	\]
    \[
	    \pause
	    \Delta x = x - x_0
	\]
	\[
		\pause
		\varepsilon\to 0 \text{ quando } \Delta x \to 0
	\]

\end{frame}

%------------------------------------------------------------------------------%
\begin{frame}{Definição -- Derivada Parcial em $x$}
  A derivada parcial, com relação a \emph{$x$},
  da função \(f\colon R\subset\R^2\to\R\), \(f(x, y)\) \\
  no ponto \((x, y) = (x_0, y_0)\)

  \[
    \frac{\partial f}{\partial x}\evalat{(x, y) = (x_0, y_0)}{}
    =\; \lim_{h\to 0} \frac{f(\emph{x_0+h}, y_0)-f(x_0, y_0)}{h}
  \]

  \vspace{\baselineskip}
  Desde que o limite exista
\end{frame}

%------------------------------------------------------------------------------%
\begin{frame}{Diferenciabilidade em $n$ Dimensões}
	A existência das derivadas parciais \emph{não} garante diferenciabilidade

	\pause
	\vspace{\baselineskip}
	Diferenciabilidade está associada a existência do \emph{plano tangente}
\end{frame}

%------------------------------------------------------------------------------%
\begin{frame}{Teorema do Incremento}
	Suponha que as derivadas parciais de primeira ordem de $f(x, y)$ sejam
	definidas em uma região aberta $R$, que contem o ponto $(x_0, y_0)$
    \pause
	e que $f_x$ e $f_y$ sejam contínuas em $(x_0, y_0)$,
	\pause
	então a variação
	\[
		\pause
        \Delta z
		= f\left(x_0+\Delta x, y_0 + \Delta y\right)
		- f(x_0, y_0)
	\]
	\pause
    com \(\left(x_0+\Delta x, y_0 + \Delta y\right)\in R\),
	\pause
    satisfaz
	\[
		\Delta z
		= f_x(x_0, y_0) \Delta x
		+ f_y(x_0, y_0) \Delta y
		+ \varepsilon_1 \Delta x
		+ \varepsilon_2 \Delta y
	\]
	\[
		\pause
        \varepsilon_1, \varepsilon_2 \to 0 \text{ quando } \Delta x, \Delta y \to 0
	\]
\end{frame}

%------------------------------------------------------------------------------%
\begin{frame}{Definição: Diferenciabilidade em 2D}
	Uma função \(z=f(x, y)\) é \emph{diferenciável} em \((x_0, y_0)\)

    \pause
    \vspace{0.5\baselineskip}
    se \(\frac{df}{dx}(x_0, y_0)\) e \(\frac{df}{dy}(x_0, y_0)\) existem

    \pause
    \vspace{0.5\baselineskip}
    e \(\Delta z = f(x_0+\Delta x, y_0+\Delta y)-f(x_o, y_0)\)

    \pause
    \vspace{0.5\baselineskip}
    satisfaz
	\[
		\Delta z
		\;=\; f_x(x_0, y_0) \Delta x
		\;+\; f_y(x_0, y_0) \Delta y
		\;+\; \varepsilon_1 \Delta x
		\;+\; \varepsilon_2 \Delta y
	\]
	\[
		\pause
        \varepsilon_1, \varepsilon_2 \to 0 \text{ quando } \Delta x, \Delta y \to 0
	\]
\end{frame}

%------------------------------------------------------------------------------%
\begin{frame}{Diferenciabilidade em 2D}
	Dizemos que $f$ é \emph{diferenciável}
	se ela é diferenciável em todos os
	pontos do seu domínio

	\pause
    \vspace{0.5\baselineskip}
    Nesse caso dizemos que seu gráfico é uma \emph{superfície lisa}
\end{frame}

%------------------------------------------------------------------------------%
\begin{frame}{Corolário do Teorema do Incremento}
    Se \(\frac{df}{dx}(x_0, y_0)\) e \(\frac{df}{dy}(x_0, y_0)\) existem

    \pause
    \vspace{0.5\baselineskip}
    e são \emph{contínuas} em uma região aberta $R$

    \pause
    \vspace{0.5\baselineskip}
    então $f$ é \emph{diferenciável} em $R$
\end{frame}

%------------------------------------------------------------------------------%
\begin{frame}{Teorema Diferenciabilidade e Contiuidade}
    Se uma função $f(x, y)$ é \emph{diferenciável} em \((x_0, y_0)\)

    \pause
    \vspace{0.5\baselineskip}
    então ela é contínua em \((x_0, y_0)\)
\end{frame}

%------------------------------------------------------------------------------%
\section{Lista Mínima}

%------------------------------------------------------------------------------%
\begin{frame}{Lista Mínima}
  Cálculo Vol. 2 do Thomas 12\textsuperscript{a} ed. -- Seção 14.3
  \begin{enumerate}
    \item Estudar o texto da seção
%    \item Resolver os exercícios:
  \end{enumerate}

  \vfill
  Atenção: A prova é baseada no livro, não nas apresentações
\end{frame}

%------------------------------------------------------------------------------%
\end{document}
%------------------------------------------------------------------------------%
